\chapter{公式与符号速查}

\section*{A.1 全局符号表}

\begin{table}[htbp]
  \centering
  \caption{全书通用符号}\label{tab:a-symbols}
  \small
  \begin{tabular}{llp{7.2cm}}
    \toprule
    符号 & 名称 & 含义 \\
    \midrule
    $W$ & 世界系 / map 系 & 全局参考坐标系，$z$ 轴向上 \\
    $O$ & odom 系 & 连续、无跳变的局部参考系 \\
    $B$ & 机体系 & 默认与 IMU 系重合 \\
    $L$、$C$ & 雷达系、相机系 & 传感器坐标系 \\
    $T_{AB}=(R_{AB},{}^{A}p_B)$ & 位姿 & 从 $B$ 系到 $A$ 系的变换 \\
    $R$、$q$ & 旋转矩阵、四元数 & $R\in\SO(3)$ \\
    $p$、$v$、$a$ & 位置、速度、加速度 & 列向量 \\
    $\omega$、$g$ & 角速度、重力 & $g=[0,0,-9.81]^{\top}$ \\
    $\tilde{\omega}$、$\tilde{a}$ & IMU 测量值 & 含零偏与噪声 \\
    $b_g$、$b_a$ & 陀螺仪/加速度计零偏 & 缓慢漂移 \\
    $x$、$\delta x$ & 状态、误差状态 & \\
    $z$、$h(\cdot)$、$r$ & 测量、观测函数、残差 & \\
    $P$、$Q$、$R$ & 状态/过程/测量协方差 & \\
    $J$、$H$ & 雅可比、海森 & \\
    $\Lambda$ & 信息矩阵 & $\Lambda=P^{-1}$ 或 $J^{\top}R^{-1}J$ \\
    $\pi=(n_\pi,d)$ & 平面参数 & $n_\pi^{\top}x=d$ \\
    \bottomrule
  \end{tabular}
\end{table}

\section*{A.2 旋转与刚体变换}

\paragraph{旋转矩阵的性质}
\begin{equation}
R^{\top}R=I,\quad \det(R)=1,\quad (R_1R_2)^{-1}=R_2^{\top}R_1^{\top}.
\end{equation}

\paragraph{Rodrigues 公式}
设 $\phi=\theta a$（$\norm{a}=1$），则
\begin{equation}
\exp(\phi^{\wedge})=I+\sin\theta\,a^{\wedge}+(1-\cos\theta)(a^{\wedge})^{2}.
\end{equation}

\paragraph{对数映射}
\begin{equation}
\theta=\arccos\frac{\tr(R)-1}{2},\qquad
\phi=\frac{\theta}{2\sin\theta}(R-R^{\top})^{\vee}.
\end{equation}

\paragraph{$\SE(3)$ 的逆}
\begin{equation}
T^{-1}=\begin{bmatrix}R^{\top}&-R^{\top}t\\0^{\top}&1\end{bmatrix}.
\end{equation}

\paragraph{位姿复合与逆}
\begin{equation}
T_{AC}=T_{AB}T_{BC},\qquad
T_{AB}^{-1}=T_{BA}.
\end{equation}

\paragraph{右扰动}
\begin{equation}
R\leftarrow R\exp(\delta\phi^{\wedge}),\qquad t\leftarrow t+R\,\delta\rho.
\end{equation}

\section*{A.3 概率与滤波}

\paragraph{贝叶斯公式}
\begin{equation}
p(x\mid z)=\frac{p(z\mid x)p(x)}{p(z)}.
\end{equation}

\paragraph{高斯分布}
\begin{equation}
p(x)=\frac{1}{(2\pi)^{n/2}\abs{\Sigma}^{1/2}}
\exp\left(-\frac12(x-\mu)^{\top}\Sigma^{-1}(x-\mu)\right).
\end{equation}

\paragraph{方差与协方差}
\begin{equation}
\Var(X)=\E[X^{2}]-(\E[X])^{2},\qquad
\Cov(X,Y)=\E[XY]-\E[X]\E[Y].
\end{equation}

\paragraph{卡尔曼滤波}
预测：
\begin{align}
\hat{x}_{k}^{-}&=F_k\hat{x}_{k-1}+B_ku_k,\\
P_k^{-}&=F_kP_{k-1}F_k^{\top}+Q_k.
\end{align}
更新：
\begin{align}
K_k&=P_k^{-}H_k^{\top}\left(H_kP_k^{-}H_k^{\top}+R_k\right)^{-1},\\
\hat{x}_k&=\hat{x}_k^{-}+K_k\left(z_k-h(\hat{x}_k^{-})\right),\\
P_k&=(I-K_kH_k)P_k^{-}.
\end{align}

\paragraph{ESKF}
误差状态传播与更新：
\begin{equation}
\delta\dot{x}=F\delta x+Gn,\qquad
K=P^{-}H^{\top}(HP^{-}H^{\top}+R)^{-1},\qquad
\delta x\leftarrow\delta x\oplus K r.
\end{equation}

\paragraph{马氏距离与卡方检验}
\begin{equation}
d^{2}=r^{\top}(HPH^{\top}+R)^{-1}r\sim\chi^{2}(k).
\end{equation}

\section*{A.4 优化}

\paragraph{最小二乘正规方程}
\begin{equation}
A^{\top}Ax=A^{\top}b.
\end{equation}

\paragraph{高斯--牛顿}
\begin{equation}
J^{\top}J\,\delta x=-J^{\top}r,\qquad x\leftarrow x+\delta x.
\end{equation}

\paragraph{Levenberg--Marquardt}
\begin{equation}
\left(J^{\top}J+\lambda I\right)\delta x=-J^{\top}r.
\end{equation}

\paragraph{Huber 鲁棒核}
\begin{equation}
\rho(s)=
\begin{cases}
s, & s\le\delta^{2},\\
2\delta\sqrt{s}-\delta^{2}, & s>\delta^{2},
\end{cases}
\qquad
w(s)=\frac{\rho'(s)}{2\sqrt{s}}=
\begin{cases}
1, & s\le\delta^{2},\\
\delta/\sqrt{s}, & s>\delta^{2}.
\end{cases}
\end{equation}

\section*{A.5 点云配准}

\paragraph{ICP 闭式旋转解}
给定去质心点对 $(p_i,q_i)$，构造
\begin{equation}
W=\sum_i p_iq_i^{\top}=U\Sigma V^{\top},\qquad
R^{*}=VU^{\top},
\end{equation}
若 $\det(R^{*})<0$，取 $V'=[v_1,v_2,-v_3]$ 后重算。平移为 $t^{*}=\bar{q}-R^{*}\bar{p}$。

\paragraph{点到面残差}
\begin{equation}
r_i=n_i^{\top}\left(Rp_i+t-q_i\right).
\end{equation}

\section*{A.6 SLAM 后端}

\paragraph{位姿图残差}
\begin{equation}
r_{ij}=\log\left(T_{ij}^{-1}T_i^{-1}T_j\right)^{\vee}\in\R^{6}.
\end{equation}

\paragraph{舒尔补}
分块系统
\begin{equation}
\begin{bmatrix}H_{cc}&H_{cp}\\H_{pc}&H_{pp}\end{bmatrix}
\begin{bmatrix}\delta x_c\\\delta x_p\end{bmatrix}
=
\begin{bmatrix}b_c\\b_p\end{bmatrix}
\end{equation}
消去相机变量后：
\begin{equation}
\left(H_{pp}-H_{pc}H_{cc}^{-1}H_{cp}\right)\delta x_p
=b_p-H_{pc}H_{cc}^{-1}b_c.
\end{equation}

\paragraph{边缘化}
用舒尔补把旧状态消去，得到新的先验因子；注意这会引入线性化误差。

\section*{A.7 导航}

\paragraph{A* 评价函数}
\begin{equation}
f(n)=g(n)+h(n),
\end{equation}
其中 $g$ 是起点到 $n$ 的实际代价，$h$ 是到目标的启发式下界（可采纳）。

\paragraph{DWA 评价函数}
\begin{equation}
G(v,\omega)=\alpha\cdot\mathrm{heading}(v,\omega)+\beta\cdot\mathrm{clearance}(v,\omega)+\gamma\cdot\mathrm{velocity}(v,\omega).
\end{equation}

\paragraph{Cartographer 扫描匹配得分}
\begin{equation}
S(T)=\sum_k M_{\mathrm{smooth}}\left(T\,h_k\right),
\end{equation}
其中 $M_{\mathrm{smooth}}$ 是双线性插值后的概率栅格。

\paragraph{控制代价（MPC）}
\begin{equation}
\min_{u_{0:N-1}}\ \sum_{k=0}^{N-1}\left(x_k-x_k^{\mathrm{ref}}\right)^{\top}Q\left(x_k-x_k^{\mathrm{ref}}\right)+u_k^{\top}Ru_k
\quad\text{s.t.}\ x_{k+1}=f(x_k,u_k).
\end{equation}
